\section{Символическая запись на исходном языке}

В настоящем разделе приводится символическая запись основных модулей программы,
определяющих состав поддерживаемых таблиц, способы генерации типовых значений и
порядок записи сформированных данных в выходные файлы.

\subsection{Модуль регистрации таблиц}

Ниже приведена символическая запись модуля, задающего состав таблиц,
поддерживаемых программой, и связывающего их с соответствующими генераторами
строк и схемами колонок.

\begin{lstlisting}[caption={Регистрация таблиц генератора}]
func register_tables():
    registry.add("region",   region_schema,   region_generator)
    registry.add("nation",   nation_schema,   nation_generator)
    registry.add("supplier", supplier_schema, supplier_generator)
    registry.add("part",     part_schema,     part_generator)
    registry.add("partsupp", partsupp_schema, partsupp_generator)
    registry.add("customer", customer_schema, customer_generator)
    registry.add("orders",   orders_schema,   orders_generator)
    registry.add("lineitem", lineitem_schema, lineitem_generator)
\end{lstlisting}

\subsection{Модуль генерации типовых значений}

Ниже приведены символические записи основных компонентов модуля генерации
типовых значений. Отдельно отражены генерация ограниченных числовых значений,
выбор токенов из распределений, построение строковых представлений и
преобразование дат TPC-H.

\begin{lstlisting}[caption={Генерация ограниченных числовых значений}]
func next_bounded_int(seed, lower_bound, upper_bound, values_per_row):
    random = RowRandomInt(seed, values_per_row)
    return random.next_int(lower_bound, upper_bound)


func next_bounded_long(seed, lower_bound, upper_bound, values_per_row):
    random = RowRandomLong(seed, values_per_row)
    return random.next_long(lower_bound, upper_bound)


func next_bounded_value(seed, lower_bound, upper_bound, scale_factor, values_per_row):
    use_64bits = scale_factor >= 30000 or upper_bound > INT32_MAX

    if use_64bits:
        return next_bounded_long(seed, lower_bound, upper_bound, values_per_row)

    return next_bounded_int(seed, lower_bound, upper_bound, values_per_row)


func advance_rows(random, row_count):
    random.advance_rows(row_count)


func finish_row(random):
    random.row_finished()
\end{lstlisting}

\begin{lstlisting}[caption={Генерация строк из распределений}]
func next_token(distribution, random):
    sample = random.next_int(0, distribution.total_weight - 1)
    cumulative = 0

    for entry in distribution:
        cumulative += entry.weight
        if sample < cumulative:
            return entry.token

    return distribution.last.token


func next_string(distribution, random):
    return next_token(distribution, random)


func next_string_sequence(distribution, count, random):
    values = copy_all_tokens(distribution)

    for current in range(0, count):
        swap_with = random.next_int(current, len(values) - 1)
        swap(values[current], values[swap_with])

    return join_with_spaces(values[0:count])
\end{lstlisting}

\begin{lstlisting}[caption={Генерация символьных и телефонных значений}]
func next_alphanumeric(average_length, random):
    length = random.next_int(0.4 * average_length, 1.6 * average_length)
    text = ""

    while len(text) < length:
        chunk = random.next_int(0, INT32_MAX)
        text += decode_base64_like_alphabet(chunk)

    return text[0:length]


func next_phone(nation_key, random):
    country_code = 10 + (nation_key mod 90)
    local1 = random.next_int(100, 999)
    local2 = random.next_int(100, 999)
    local3 = random.next_int(1000, 9999)
    return format("%02d-%03d-%03d-%04d", country_code, local1, local2, local3)
\end{lstlisting}

\begin{lstlisting}[caption={Преобразование дат TPC-H}]
func to_unix_epoch(generated_date):
    return (generated_date - MIN_GENERATE_DATE) + UNIX_EPOCH_OFFSET


func is_in_past(generated_date):
    return to_julian(generated_date) <= CURRENT_DATE


func format_tpch_date(generated_date):
    year, month, day = to_ymd(generated_date)
    return format("%04d-%02d-%02d", year, month, day)
\end{lstlisting}

\subsection{Модуль формирования пула текста}

Ниже приведена символическая запись модуля \texttt{TextPool}, используемого для
построения текстовых атрибутов. Пул заполняется синтетическими предложениями,
построенными по грамматическим шаблонам TPC-H, после чего генераторы таблиц
через компонент \texttt{RandomText} берут из него подстроки без дополнительного
конструирования текста для каждой строки.

\begin{lstlisting}[caption={Формирование и использование TextPool}]
func build_text_pool(target_size):
    random = RowRandomInt(933588178, MAX_INT)
    text = ""

    while len(text) < target_size:
        text += generate_sentence(random)

    return text[0:target_size]


func generate_sentence(random):
    sentence = ""
    syntax = next_token(grammar_distribution, random)

    for symbol in syntax:
        if symbol == "V":
            sentence += generate_verb_phrase(random)
        elif symbol == "N":
            sentence += generate_noun_phrase(random)
        elif symbol == "P":
            sentence += next_token(prepositions, random) + " the "
            sentence += generate_noun_phrase(random)
        elif symbol == "T":
            sentence = trim_trailing_space(sentence)
            sentence += next_token(terminators, random)

        sentence += " "

    return sentence


func next_text(text_pool, average_length, random):
    min_length = 0.4 * average_length
    max_length = 1.6 * average_length

    offset = random.next_int(0, text_pool.size - max_length)
    length = random.next_int(min_length, max_length)

    /* (*@\texttt{Возврат представления}@*) */
    return text_pool.view(offset, offset + length)
\end{lstlisting}

\subsection{Модуль выбора стратегии записи}

Ниже приведена символическая запись модуля, определяющего способ записи
сформированных данных. Для форматов \texttt{Parquet} и \texttt{ORC}
поддерживаются стратегии \texttt{parallel-files} и
\texttt{single-file-ordered}; для форматов \texttt{Lance} и \texttt{Vortex}
используется только упорядоченная запись в один логический набор файлов.

\begin{lstlisting}[caption={Выбор стратегии записи}]
func generate_table_with_strategy(context, table, output, writer_options, format_driver, scheduler_options):
    part_count = format_driver.resolve_part_count(table, context.scale, writer_options)
    strategy = resolve_selected_write_strategy(scheduler_options, format_driver)

    if strategy == "parallel-files":
        return run_parallel_partition_files(
            context, table, output, writer_options, format_driver, part_count
        )

    if format_driver.ordered_execution_model() == "native-multiplexed":
        return run_native_multiplexed_ordered(
            context, table, output, writer_options, format_driver, part_count
        )

    return run_foreign_streaming_ordered(
        context, table, output, writer_options, format_driver, part_count
    )


func resolve_selected_write_strategy(options, format_driver):
    if options.write_strategy == "auto":
        strategy = format_driver.preferred_strategy()
    else:
        strategy = options.write_strategy

    assert format_driver.supports_strategy(strategy)
    return strategy
\end{lstlisting}

\subsection{Модуль записи в формат Parquet}

Ниже приведена символическая запись модуля записи в формат \texttt{Parquet}. В
зависимости от выбранной стратегии данные либо распределяются по нескольким
файлам партиций, либо несколькими производителями готовятся параллельно и затем
в порядке номеров партиций объединяются в один файл
\texttt{table-1.parquet}.

\begin{lstlisting}[caption={Запись таблицы в формат Parquet}]
func parquet_parallel_partition_files(table, partition, batches, output, options):
    path = output / table.name / format_parquet_partition_name(table, partition)
    temp_path = path + ".tmp"

    sink = open_output_stream(temp_path)
    sink = maybe_buffer(sink, options.output_buffer_bytes)

    writer = open_parquet_writer(
        schema = table.schema,
        compression = options.compression,
        row_group_rows = resolve_row_group_rows(table, options),
        use_threads = options.use_threads
    )

    for batch in batches:
        writer.write_record_batch(batch)

    writer.close()
    sink.close()
    move_file(temp_path, path)


func parquet_single_file_ordered(table, output, options, part_count):
    sink = open_single_output(output / table.name / (table.name + "-1.parquet"))
    producers = spawn_partition_generators(part_count)
    states = {}
    expected_part = 1

    while message = read_batch_message(producers):
        states[message.part].append(message)

        while states[expected_part].is_ready():
            /* (*@\texttt{Партиция равна row group}@*) */
            sink.begin_partition()

            for batch in states[expected_part].batches:
                sink.write_batch(batch)

            sink.end_partition()
            expected_part += 1

    sink.close()
\end{lstlisting}

\subsection{Модуль записи в формат ORC}

Ниже приведена символическая запись модуля записи в формат \texttt{ORC}. При
стратегии \texttt{parallel-files} каждая партиция записывается в собственный
файл, а при стратегии \texttt{single-file-ordered} батчи партиций
последовательно направляются в единый ORC writer с сохранением порядка
партиций.

\begin{lstlisting}[caption={Запись таблицы в формат ORC}]
func orc_parallel_partition_files(table, partition, batches, output, options):
    path = output / table.name / format_orc_partition_name(table, partition)
    temp_path = path + ".tmp"

    sink = open_output_stream(temp_path)
    sink = maybe_buffer(sink, options.output_buffer_bytes)

    writer = open_orc_writer(
        schema = table.schema,
        compression = options.compression,
        stripe_bytes = resolve_stripe_bytes(table, options)
    )

    for batch in batches:
        writer.write(batch)

    writer.close()
    sink.close()
    move_file(temp_path, path)


func orc_single_file_ordered(table, output, options, part_count):
    sink = open_single_output(output / table.name / (table.name + "-1.orc"))
    producers = spawn_partition_generators(part_count)
    states = {}
    expected_part = 1

    while message = read_batch_message(producers):
        states[message.part].append(message)

        while states[expected_part].is_ready():
            /* (*@\texttt{Явного начала stripe нет}@*) */
            sink.begin_partition()

            for batch in states[expected_part].batches:
                sink.write_batch(batch)

            sink.end_partition()
            expected_part += 1

    sink.close()
\end{lstlisting}

\subsection{Модуль записи в формат Lance}

Ниже приведена символическая запись модуля записи в формат \texttt{Lance}. Для
данного формата используется только стратегия
\texttt{single-file-ordered}: генератор последовательно подает партиции в
упорядоченный writer, а фактическая запись набора данных выполняется через
мост в реализацию на Rust.

\begin{lstlisting}[caption={Запись таблицы в формат Lance}]
func lance_single_file_ordered(table, output, options, part_count):
    dataset_path = output / table.name / (table.name + "-1.lance")
    schema = export_arrow_schema(table.schema)
    bridge_options = build_lance_bridge_options(options)

    /* (*@\texttt{Вызов Rust FFI}@*) */
    handle = lance_writer_open(dataset_path, schema, bridge_options)

    for part in range(1, part_count + 1):
        /* (*@\texttt{Вызов Rust FFI}@*) */
        lance_writer_begin_partition(handle, part, part_count)

        for batch in generate_partition_batches(table, part):
            if batch.row_count == 0:
                continue

            arrow_batch = export_record_batch(batch)

            /* (*@\texttt{Вызов Rust FFI}@*) */
            lance_writer_push_batch(handle, arrow_batch)

        /* (*@\texttt{Вызов Rust FFI}@*) */
        lance_writer_end_partition(handle)

    /* (*@\texttt{Вызов Rust FFI}@*) */
    lance_writer_finish(handle)
    lance_writer_destroy(handle)
\end{lstlisting}

\subsection{Модуль записи в формат Vortex}

Ниже приведена символическая запись модуля записи в формат \texttt{Vortex}. Как
и в случае \texttt{Lance}, запись выполняется только в режиме
\texttt{single-file-ordered}; сериализация батчей и финализация файла
осуществляются через мост в реализацию на Rust.

\begin{lstlisting}[caption={Запись таблицы в формат Vortex}]
func vortex_single_file_ordered(table, output, options, part_count):
    file_path = output / table.name / (table.name + "-1.vortex")
    schema = export_arrow_schema(table.schema)
    bridge_options = build_vortex_bridge_options(options)
    wrote_non_empty_batch = false

    /* (*@\texttt{Вызов Rust FFI}@*) */
    handle = vortex_writer_open(file_path, schema, bridge_options)

    for part in range(1, part_count + 1):
        /* (*@\texttt{Вызов Rust FFI}@*) */
        vortex_writer_begin_partition(handle, part, part_count)

        for batch in generate_partition_batches(table, part):
            if batch.row_count == 0:
                continue

            arrow_batch = export_record_batch(batch)

            /* (*@\texttt{Вызов Rust FFI}@*) */
            vortex_writer_push_batch(handle, arrow_batch)
            wrote_non_empty_batch = true

        /* (*@\texttt{Вызов Rust FFI}@*) */
        vortex_writer_end_partition(handle)

    if not wrote_non_empty_batch:
        empty_batch = make_empty_record_batch(table.schema)

        /* (*@\texttt{Вызов Rust FFI}@*) */
        vortex_writer_push_batch(handle, export_record_batch(empty_batch))

    /* (*@\texttt{Вызов Rust FFI}@*) */
    vortex_writer_finish(handle)
    vortex_writer_destroy(handle)
\end{lstlisting}
